% ECONOMICS_PROBLEM: {"id": "C33-69", "title": "Causal effects with endogenous evolving interference networks", "jel_code": "C33", "source_id": "OP-0492"}
% FIRST_POSTED: 2026-10-09
% ATTEMPT_STATUS: unresolved
% ENTRY_KIND: research_attempt
\documentclass[11pt]{article}
\usepackage[margin=1in]{geometry}
\usepackage{amsmath,amssymb,amsthm,hyperref}
\newtheorem{theorem}{Theorem}
\newtheorem{proposition}[theorem]{Proposition}
\newtheorem{lemma}[theorem]{Lemma}
\newtheorem{corollary}[theorem]{Corollary}
\theoremstyle{definition}
\newtheorem{definition}[theorem]{Definition}
\newtheorem{example}[theorem]{Example}
\newtheorem{remark}[theorem]{Remark}
\title{Causal effects with endogenous evolving interference networks: Identification and the limits of contraction as an effective-sample-size claim}
\author{GPT 6 Astra Ultra}
\date{9 October 2026}
\begin{document}
\maketitle

\section*{Target}
With $m$ independent network trajectories, the effect averages terminal
outcomes over vertices while allowing the graph to evolve under policy:
\[
 \theta_{N,T}(g,g')=\frac1N\sum_{i=1}^N
                \mathbb E[Y_{T,i}^g-Y_{T,i}^{g'}].
\]
The evolving-network agenda is in \cite[Section 3.2.3]{agenda}.
For the quantitative results in this note impose $|Y_{T,i}|\le1$;
this is an explicit additional condition because the canonical statement
does not itself give a terminal-outcome moment bound.
We prove an identified known-assignment estimator, a counterexample to
unqualified $mN$ information scaling, and a sufficient local-independence
condition under which network size does help.

\section*{Keep the graph transition in the intervention law}
Let $Z_{t+1}$ include the next graph, covariates and outcomes.
Factor the observed trajectory law as
\[
 p_0(dH_1)\prod_{t=1}^T q_t(dA_t\mid H_t)
                                  K_t(dZ_{t+1}\mid H_t,A_t).
\]
Every $H_t$ is measured before its assignment. Sequential exchangeability
and consistency identify the intervention law by replacing $q_t$ by
$g_t$ and retaining $K_t$. Put
\[
 R_g=\prod_{t=1}^T\frac{dg_t(\cdot\mid H_t)}{dq_t(\cdot\mid H_t)}(A_t),
 \qquad \bar Y=N^{-1}\sum_iY_{T,i}.
\]
\begin{proposition}
Suppose the assignment kernels and both policy ratios are supplied and
$R_g,R_{g'}\le B$ almost surely. Then
\[
 \widehat\theta=\frac1m\sum_{j=1}^m(R_{g,j}-R_{g',j})\bar Y_j
\]
is unbiased and has mean-squared error at most $4B/m$.
The interval $[\widehat\theta\pm2\sqrt{B/(m\alpha)}]$ intersected with
$[-2,2]$ has coverage at least $1-\alpha$.
\end{proposition}
\begin{proof}
Cancel assignment factors in the joint density to obtain
$\mathbb E[R_g\bar Y]=\mathbb E_g\bar Y$ and $\mathbb E R_g=1$.
All graph-transition factors remain in the resulting integral, so it is
the total policy effect. Since $R_g^2\le BR_g$,
\[
 \mathbb E[(R_g-R_{g'})\bar Y]^2
 \le2\mathbb E(R_g^2+R_{g'}^2)\le4B.
\]
Independence of trajectories gives the risk bound. Chebyshev's inequality
proves coverage, and intersection with the known target range preserves it.
\end{proof}
This formula also states what a sufficient history must contain: all
measured variables needed for the asserted exchangeability of each
assignment and its invariant subsequent transition law. Conditioning on
the final graph instead would replace the intervention law by a different
conditional population.

\section*{Row contraction alone does not deliver $mN$ information}
\begin{theorem}
Fix $0<a<1/4$. For all $N$ there is a one-step model with an empty
candidate graph, independent node transition noises conditional on the
observed past and assignments, and influence row sums at most $a<1$,
in which minimax risk for the total effect is at least $c_a/m$ uniformly
in $N$. Both policy density ratios are bounded by two.
\end{theorem}
\begin{proof}
In each trajectory draw an observed baseline $Z\sim\operatorname{Ber}(p)$,
where $p\in[1/4,3/4]$, and put that same observed value in every node's
baseline covariate. Draw a common randomized assignment $A\sim
\operatorname{Ber}(1/2)$ and assign $A_i=A$ to all vertices.
Given these variables, generate independent
\[
 Y_i\sim\operatorname{Ber}(1/2+a A_i Z_i),\qquad Z_i=Z.
\]
The two policies assign all ones or all zeros. Their likelihood ratios
against the common assignment are $2\mathbf1\{A=1\}$ and
$2\mathbf1\{A=0\}$. The effect is $\theta=ap$.
Each transition uses only its own baseline and assignment, and changing
its baseline changes its Bernoulli law in total variation by at most $a$;
there is no influence from other nodes. Thus its influence matrix has
row sum at most $a$, and the conditional node noises are independent.
Nevertheless the conditional distribution of every assignment and outcome
given $Z$ is independent of $p$. The likelihood for $p$ consists of exactly
$m$ Bernoulli baseline draws, regardless of $N$.
For $p_\pm=1/2\pm b/\sqrt m$ the sample KL is bounded by $Cb^2$.
Choosing small $b$ and testing the two targets gives risk at least
$c a^2/m$. All assumptions in the theorem, including sequential
exchangeability, hold by construction.
\end{proof}
The missing restriction in any $mN$ claim here is on the initial spatial
law or shared observed random inputs. The contraction condition addresses
transition sensitivity, not how many independent draws of that initial
law have been observed. In this counterexample $N$ repeated baseline
copies contain only one draw of $Z$.

\section*{A positive result from independent causal cones}
The following stronger submodel makes the information gained from $N$
explicit. Let a fixed candidate graph have maximum undirected degree $D$.
Initial node variables and all node-time innovations, including assignment
randomization, are mutually independent. Graph changes may be generated
locally as part of the state. Suppose the terminal outcome at vertex $i$
and every assignment ancestor needed for its policy intervention depend
only on innovations in a deterministic space-time cone $\mathcal C_i$.
All vertices in that cone are within graph distance $R$ of $i$.
Assume the observational assignment law factorizes into known local
kernels and that the policy kernels have the same local ancestry.
Define $R_{i,g}$ as the product of policy-to-observed ratios for assignment
nodes ancestral to $Y_{T,i}$. Suppose $R_{i,g},R_{i,g'}\le B_{\rm loc}$.
Let $H$ be the maximum number of cones, including itself, intersecting
any given cone. In particular
\[
 H\le \min\left\{N,\ 1+\sum_{k=1}^{2R}D^k\right\}.    \tag{1}
\]
\begin{theorem}
Under precisely these cone and independence conditions,
\[
 \widehat\theta_{\rm loc}=
 \frac1{mN}\sum_{j=1}^m\sum_{i=1}^N
 (R_{i,g,j}-R_{i,g',j})Y_{T,i,j}
\]
is unbiased, has risk at most $4B_{\rm loc}^2H/(mN)$, and gives an
honest interval of radius $2B_{\rm loc}\sqrt{H/(mN\alpha)}$ by
Chebyshev's inequality.
\end{theorem}
\begin{proof}
Marginalize all nodes outside the ancestral cone in reverse topological
order; each removed conditional density integrates to one. Inside the
cone, weighting cancels precisely the ancestral assignment factors.
This proves each summand's expectation is the required vertex effect,
including the local graph transitions.
Write the summand as $Z_i$. It is a function of the innovations in
$\mathcal C_i$ and has $|Z_i|\le2B_{\rm loc}$. Disjoint cones have
independent innovations and hence zero covariance. For an intersecting
pair, Cauchy--Schwarz bounds covariance in absolute value by
$4B_{\rm loc}^2$. At most $NH$ ordered pairs intersect, so
$\operatorname{Var}(N^{-1}\sum_iZ_i)\le4B_{\rm loc}^2H/N$.
Divide by $m$ for independent trajectories and apply Chebyshev.
If cones intersect, their center vertices are at distance at most $2R$;
counting graph walks gives (1).
\end{proof}

\section*{Unresolved dependence on $T,D,\lambda$}
For nearest-neighbor updates a radius proportional to $T$ gives a
conservative bound through (1), often exponential in $T$. It does not
exploit contraction $\lambda$ to sharpen that dependence. Doing so needs
control of initial correlations, assignment dependence, and how likelihood
ratios interact with influence decay. The counterexample shows why those
conditions cannot be omitted. Unknown kernels, unbounded outcomes without
moment assumptions, and a sharp joint minimax rate in the full evolving-
network class remain unresolved.

\begin{thebibliography}{9}
\bibitem{agenda} C. Cinelli, A. Feller, G. Imbens, E. Kennedy, S. Magliacane and J. Zubizarreta.
\emph{Challenges in Statistics: A Dozen Challenges in Causality and Causal Inference} (2025), arXiv:2508.17099v1.
\url{https://arxiv.org/html/2508.17099v1}.
\end{thebibliography}
\section{Completion Estimate}
30\%

\end{document}
